\documentclass[11pt]{article}

\usepackage[T1]{fontenc}
\usepackage{amsmath,amssymb,amsthm,mathtools}
\usepackage[margin=0.82in]{geometry}
\usepackage{booktabs,longtable,array}

\newtheorem{lemma}{Lemma}
\newtheorem{proposition}[lemma]{Proposition}
\newtheorem{corollary}[lemma]{Corollary}

\newcommand{\W}{\mathcal W}
\newcommand{\Dlt}{\Delta}
\newcommand{\pimin}{\pi_-}
\newcommand{\pimax}{\pi_+}
\newcolumntype{L}[1]{>{\raggedright\arraybackslash}p{#1}}

\title{Connected High-Event Registry}
\author{}
\date{Analytic ledger replacement, Packet C}

\begin{document}
\maketitle

\begin{abstract}
This note replaces the exact-executable arithmetic used in the high
spherical-shell staircase by displayed rational witnesses.  It contains
the complete seventeen-face localization registry, the three algebraic
splits, the three dashed (impossible) cap cells, the five global checkpoint
payments, the seven conditional payments, and the corrected cap--74 path.
It also contains the order-free closed-post-event theorem which maps every
high-staircase state to a payment at its own event wall.  That theorem,
rather than the historical executable row census, is the finite-combinatorial
closure used by the proof.
All comparisons are strict rational comparisons after the two elementary
bounds \(333/106<\pi<22/7\).  In particular, no decimal approximation,
interval package, or external ledger is a logical input to this packet.
\end{abstract}

\section{Definitions and strict proxy convention}

For \(0<\rho<1\), put
\begin{equation}
 z=z_\rho:=\frac{\pi}{1-\rho},\qquad
 x=x_\rho:=z_\rho^2,\qquad
 k_m(\rho):=\sqrt{x_\rho+m(m+1)}.
 \label{eq:def-z-k}
\end{equation}
The three-dimensional Weyl term is
\begin{equation}
 \W(\rho,K):=\frac{2}{9\pi}(1-\rho^3)K^3.
 \label{eq:def-W}
\end{equation}
At \(K=k_m(\rho)\), the product part of the strict channel proxy for the
mode \((\ell,n)\), \(n\geq2\), is
\[
 n^2x+\ell(\ell+1)<x+m(m+1).
\]
It is therefore governed by
\begin{equation}
 \Dlt_{m;n,\ell}(\rho)
 :=m(m+1)-\ell(\ell+1)-(n^2-1)x_\rho.
 \label{eq:delta}
\end{equation}
The channel is product-possible only when \(\Dlt_{m;n,\ell}>0\).
Equality is excluded because the spectral proxy is strict.  Since
\begin{equation}
 x_\rho'=\frac{2\pi^2}{(1-\rho)^3}>0,
 \qquad
 \Dlt_{m;n,\ell}'=-(n^2-1)x_\rho'<0,
 \label{eq:delta-monotone}
\end{equation}
a positive witness at \(r\) controls the side \(\rho\leq r\), while a
negative witness controls the side \(\rho\geq r\).

We write \(H\) for the largest first-mode angular index retained in a cap
cell and \(h\) for the largest second-mode angular index.  A complete block
through \(H\) has multiplicity \((H+1)^2\), because
\(\sum_{\ell=0}^{H}(2\ell+1)=(H+1)^2\).

\section{The only transcendental enclosure}

\begin{lemma}[Rational enclosure of \(\pi\)]
\label{lem:pi}
One has
\begin{equation}
 \pimin:=\frac{333}{106}<\pi<\frac{22}{7}=: \pimax.
 \label{eq:pi-bounds}
\end{equation}
\end{lemma}

\begin{proof}
For \(0<u<1\), the alternating expansion of \(\arctan u\) lies between
each pair of consecutive partial sums.  The tangent addition formula gives
\[
 \tan\!\left(4\arctan\frac15-\arctan\frac1{239}\right)=1;
\]
indeed, \(\tan(4\arctan(1/5))=120/119\), so both the numerator and
denominator in the last tangent quotient are \(28561/(119\cdot239)\).
The angle lies in \((0,\pi/2)\), and hence Machin's identity follows:
\[
 \pi=16\arctan\frac15-4\arctan\frac1{239}.
\]
Using an even partial sum for the positive term and an upper partial sum
for the subtracted term gives
\begin{align*}
 \pi
 &>16\left(\frac15-\frac1{3\cdot5^3}
       +\frac1{5\cdot5^5}-\frac1{7\cdot5^7}\right)-\frac4{239},\\
 \left[\text{right side}\right]-\frac{333}{106}
 &=\frac{3418213}{41563593750}>0.
\end{align*}
Likewise,
\begin{align*}
 \pi
 &<16\left(\frac15-\frac1{3\cdot5^3}
       +\frac1{5\cdot5^5}\right)
       -4\left(\frac1{239}-\frac1{3\cdot239^3}\right),\\
 \frac{22}{7}-\left[\text{right side}\right]
 &=\frac{1845738322}{1493178640625}>0.
\end{align*}
This proves \eqref{eq:pi-bounds}.
\end{proof}

\section{The seventeen rational localization faces}

Let
\[
 \rho_c:=\frac1{1+2\pi}.
\]
Because \(\pi>3\), \(\rho_c<1/7\).  The seventeen registry ratios are
strictly ordered as follows:
\begin{equation}
\begin{split}
 \frac4{25}<\frac15<\frac14<\frac3{10}<\frac13<\frac7{20}
 <\frac38<\frac25<\frac5{12}<\frac37<\frac12\\
 <\frac8{15}<\frac{107}{200}<\frac35<\frac{16}{25}
 <\frac{13}{20}<\frac23.
\end{split}
\label{eq:ratio-order}
\end{equation}
For a completely explicit distinctness check, the sixteen consecutive
gaps in \eqref{eq:ratio-order} are
\begin{equation}
 \frac1{25},\frac1{20},\frac1{20},\frac1{30},\frac1{60},
 \frac1{40},\frac1{40},\frac1{60},\frac1{84},\frac1{14},
 \frac1{30},\frac1{600},\frac{13}{200},\frac1{25},
 \frac1{100},\frac1{60}.
 \label{eq:ratio-gaps}
\end{equation}
Moreover
\[
 \frac4{25}-\frac17=\frac3{175}>0,
 \qquad
 \frac78-\frac23=\frac5{24}>0.
\]
Thus all seventeen ratios lie in the high strip
\(\rho_c<\rho<7/8\).

For a positive row below, we use
\[
 \Dlt_{m;n,\ell}(r)>
 m(m+1)-\ell(\ell+1)
 -(n^2-1)\left(\frac{\pimax}{1-r}\right)^2.
\]
For a negative row, we use
\[
 -\Dlt_{m;n,\ell}(r)>
 (n^2-1)\left(\frac{\pimin}{1-r}\right)^2
 -m(m+1)+\ell(\ell+1).
\]
The resulting exact witnesses are the following.  ``Retain'' means only
that the product wall does not permit deletion of the indicated block;
the independent Bessel wall may still remove it.

\begingroup
\small
\renewcommand{\arraystretch}{1.16}
\begin{longtable}{@{}c c c L{3.15cm} L{5.35cm}@{}}
\toprule
no.&\(r\)&\((m;n,\ell)\)&exact sign witness&one-sided cap/event consequence\\
\midrule
\endfirsthead
\toprule
no.&\(r\)&\((m;n,\ell)\)&exact sign witness&one-sided cap/event consequence\\
\midrule
\endhead
1&\(1/5\)&\((11;3,1)\)&
 \(\Dlt>320/49\)&retain the third block \(0{:}1\) for \(\rho\leq1/5\)\\
2&\(3/10\)&\((8;2,3)\)&
 \(-\Dlt>58215/137641\)&at \(k_8\), \(h\leq2\) for \(\rho\geq3/10\)\\
3&\(1/3\)&\((8;2,2)\)&
 \(-\Dlt>27699/44944\)&at \(k_8\), \(h\leq1\) for \(\rho\geq1/3\)\\
4&\(3/8\)&\((9;2,3)\)&
 \(\Dlt>2622/1225\)&retain the \(h=3\) column below this face\\
5&\(2/5\)&\((9;2,2)\)&
 \(\Dlt>248/147\)&retain the \(h=2\) column below this face\\
6&\(5/12\)&\((9;2,1)\)&
 \(\Dlt>2200/2401\)&retain the \(h=1\) column below this face\\
7&\(3/7\)&\((9;2,0)\)&
 \(-\Dlt>120843/179776\)&no second mode at \(k_9\) on the upper side\\
8&\(1/2\)&\((10;2,0)\)&
 \(-\Dlt>23677/2809\)&no second mode at \(k_{10}\) on the upper side\\
9&\(107/200\)&\((11;2,0)\)&
 \(-\Dlt>13302732/2699449\)&no second mode at \(k_{11}\) on the upper side\\
10&\(8/15\)&\((11;2,0)\)&
 \(-\Dlt>2175627/550564\)&the clean stated cutoff: a second mode forces \(\rho<8/15\)\\
11&\(3/5\)&\((8;2,0)\)&
 \(-\Dlt>5080707/44944\)&select the no-second \(k_8\) cap/payment side\\
12&\(16/25\)&\((9;2,0)\)&
 \(-\Dlt>1555635/11236\)&select the no-second \(k_9\) cap/payment side\\
13&\(13/20\)&\((10;2,0)\)&
 \(-\Dlt>18126190/137641\)&select the no-second \(k_{10}\) cap/payment side\\
14&\(2/3\)&\((11;2,0)\)&
 \(-\Dlt>1510851/11236\)&select the no-second \(k_{11}\) cap/payment side\\
15&\(4/25\)&\((10;3,0)\)&
 \(-\Dlt>260740/137641\)&no third mode at \(k_{10}\) on the upper side\\
16&\(1/4\)&\((11;3,0)\)&
 \(-\Dlt>23484/2809\)&a third mode at \(k_{11}\) forces \(\rho<1/4\)\\
17&\(7/20\)&\((9;2,4)\)&
 \(-\Dlt>36230/474721\)&the \(H=6,h=4\) cap--74 cell can occur only below this face\\
\bottomrule
\end{longtable}
\endgroup

Every entry is a displayed positive rational number.  Together with the
monotonicity \eqref{eq:delta-monotone}, the table proves the asserted side
on an interval, not merely at a sampled point.

\section{Three algebraic splits and the five \(k_{11}\) implications}

At \(\rho=\rho_c\), \(x=(\pi+1/2)^2\), while at \(\rho=7/8\),
\(x=64\pi^2\).  Lemma~\ref{lem:pi} gives the exact chain
\begin{equation}
 (\pi+\tfrac12)^2
 <\left(\frac{51}{14}\right)^2
 =\frac{2601}{196}<16<\frac{68}{3}<34<64\cdot3^2<64\pi^2.
 \label{eq:split-strip}
\end{equation}
Consequently each of \(x=16,68/3,34\) is attained at a unique point in
the high strip.  The relevant channel arithmetic is
\begin{equation}
\begin{array}{c|c|c}
x&\text{channel}&\text{exact value of }\Dlt\\ \hline
16&(m;n,\ell)=(11;3,1)&132-2-8\cdot16=2>0,\\
68/3&(10;2,6)&110-42-3(68/3)=0,\\
34&(11;2,5)&132-30-3\cdot34=0.
\end{array}
\label{eq:algebraic-splits}
\end{equation}
Thus the first split lies on the possible side.  The latter two are exact
channel walls, and their new modes are excluded by the strict convention.

The two localization consequences at \(k_{11}\) may also be read without
the table:
\begin{align}
 x_{8/15}-44
 &>\left(\frac{\pimin}{1-8/15}\right)^2-44
 =\frac{725209}{550564}>0,\label{eq:k11-second}\\
 x_{1/4}-\frac{33}{2}
 &>\left(\frac{\pimin}{1-1/4}\right)^2-\frac{33}{2}
 =\frac{5871}{5618}>0.\label{eq:k11-third}
\end{align}
Hence a second mode at \(k_{11}\) forces \(\rho<8/15\), and a third mode
forces \(\rho<1/4\).  Finally, the first-mode ball walls give
\begin{align}
 \left(\frac{69}{5}\right)^2-132
 &=\frac{1461}{25}
 =44+\frac{361}{25}
 =\frac{33}{2}+\frac{2097}{50},\label{eq:H10}\\
 \left(\frac{64}{5}\right)^2-132
 &=\frac{796}{25}
 =\frac{33}{2}+\frac{767}{50}.\label{eq:H9}
\end{align}
If \(H=10\) is present at \(k_{11}\), then \(x\) is larger than both the
second threshold \(44\) and the third threshold \(33/2\); hence every
second and third mode is excluded.  If \(H=9\) is present, every third
mode is excluded.  Equations \eqref{eq:k11-second}--\eqref{eq:H9} are the
five formerly separate \(k_{11}\) localization implications.

\section{The three impossible cap cells}

If the first mode \(H\) is counted at \(k_m\), its strict ball wall
\(\beta_{H,1}<k_m\) forces
\(x>\beta_{H,1}^2-m(m+1)\).  If the second mode \(h\) is counted, its
product wall forces
\(x<(m(m+1)-h(h+1))/3\).  The three dashed cells are therefore empty by
the following exact, incompatible bounds:
\begin{equation}
\begin{array}{c|c|c|c}
\text{cell}&\text{required lower bound}&\text{required upper bound}
&\text{gap}\\ \hline
k_8: H=6,\ h\ge2&100-72=28&(72-6)/3=22&6\\
k_9: H=8,\ h\ge0&(71/6)^2-90=1801/36&90/3=30&721/36\\
k_9: H=7,\ h\ge3&(23/2)^2-90=169/4&(90-12)/3=26&65/4.
\end{array}
\label{eq:impossible-cells}
\end{equation}
No limiting equality can rescue a cell: both necessary modal inequalities
are strict.

\section{Checkpoint monotonicity and five global payments}

Put \(M=m(m+1)\) and
\(\W_m(\rho):=\W(\rho,k_m(\rho))\).  Logarithmic differentiation shows
that \(\W_m'\) has the sign of
\begin{equation}
 \pi^2(1+\rho)-M\rho^2(1-\rho)^2.
 \label{eq:Wm-derivative}
\end{equation}
For \(m\leq11\),
\[
 \pi^2(1+\rho)-M\rho^2(1-\rho)^2
 >9-\frac{132}{16}=\frac34>0.
\]
Thus every \(\W_m\), \(m\leq11\), is strictly increasing.  This single
argument replaces both executable derivative-bound predicates.

At \(\rho_c\), one has \(z=\pi+1/2>193/53\),
\(\rho_c<1/7\), and \(1/\pi>7/22\).  Consequently
\begin{equation}
 \W(\rho_c,K)>
 \frac{2}{9}\frac7{22}\left(1-\frac1{7^3}\right)K^3
 =\frac{38}{539}K^3.
 \label{eq:global-coarse}
\end{equation}
For each row in the next table, the middle column proves
\(k_m(\rho_c)>t_m\), and the last column proves
\(\W_m(\rho_c)>C_m\).

\begin{center}
\small
\renewcommand{\arraystretch}{1.18}
\begin{tabular}{@{}c c c c c@{}}
\toprule
\(m\)&cap \(C_m\)&\(t_m\)&
\((193/53)^2+M-t_m^2\)&\((38/539)t_m^3-C_m\)\\
\midrule
7&40&\(208/25\)&\(67049/1755625\)&\(5083656/8421875\)\\
8&53&\(923/100\)&\(1901439/28090000\)&\(656778873/269500000\)\\
9&73&\(254/25\)&\(61431/1755625\)&\(7911557/8421875\)\\
10&89&\(111/10\)&\(14211/280900\)&\(1999489/269500\)\\
11&121&\(241/20\)&\(65271/1123600\)&\(5076899/2156000\)\\
\bottomrule
\end{tabular}
\end{center}
Every last-column entry is positive.  Monotonicity therefore extends each
global payment from \(\rho_c\) to the complete high strip.

\section{Seven conditional first-block payments}

We now give a general reduction which avoids any numerical enclosure of a
radical.  Suppose the presence of first mode \(H\) below the checkpoint
\(k_m\) forces \(K>t\).  Set
\begin{equation}
 a^2:=t^2-m(m+1),\qquad
 \rho_0:=1-\frac{\pi}{a}.
 \label{eq:conditional-rho0}
\end{equation}
All seven rows below have \(a>\pi\): the smallest listed radicand is
\[
 \frac{89}{4},\qquad
 \frac{89}{4}-\left(\frac{22}{7}\right)^2
 =\frac{2425}{196}>0.
\]
If the mode is present for some
\(K\leq k_m(\rho)\), then \(k_m(\rho)>t\), hence
\(x_\rho>a^2\) and \(\rho>\rho_0\).  By the strict increase of
\(\W_m\),
\[
 \W_m(\rho)>\W_m(\rho_0)=\W(\rho_0,t).
\]
Choose a rational \(r>\rho_0\).  Since \(\W(\rho,t)\) decreases with
\(\rho\) at fixed \(t\), Lemma~\ref{lem:pi} yields
\begin{equation}
 \W(\rho_0,t)>\W(r,t)>
 \frac7{99}(1-r^3)t^3.
 \label{eq:conditional-coarse}
\end{equation}
The condition \(r>\rho_0\) follows, without taking a square root, from
\begin{equation}
 L:=\pimin^2-(1-r)^2a^2>0.
 \label{eq:location-witness}
\end{equation}
Here is the complete calculation.  The column \(P\) is the reserve in
\eqref{eq:conditional-coarse} above the stated cap.

\begingroup
\small
\renewcommand{\arraystretch}{1.18}
\begin{longtable}{@{}c c c c c c c c@{}}
\toprule
\(m\)&\(H\)&\(t\)&cap&\(a^2\)&\(r\)&\(L\)&
\(P=\frac7{99}(1-r^3)t^3-\mathrm{cap}\)\\
\midrule
\endfirsthead
\toprule
\(m\)&\(H\)&\(t\)&cap&\(a^2\)&\(r\)&\(L\)&\(P\)\\
\midrule
\endhead
7&6&\(10\)&53&44&\(8/15\)&\(725209/2528100\)&\(18659/2673\)\\
8&7&\(23/2\)&80&\(241/4\)&\(3/5\)&\(64349/280900\)&\(213281/49500\)\\
9&8&\(71/6\)&81&\(1801/36\)&\(3/5\)&\(4713989/2528100\)&\(14506973/1336500\)\\
10&7&\(23/2\)&100&\(89/4\)&\(7/20\)&\(2105431/4494400\)&\(18539033/6336000\)\\
10&8&\(71/6\)&97&\(1081/36\)&\(3/7\)&\(317585/4955076\)&\(2865431/261954\)\\
10&9&\(64/5\)&100&\(1346/25\)&\(3/5\)&\(8811001/7022500\)&\(25143284/1546875\)\\
11&9&\(64/5\)&125&\(796/25\)&\(1/2\)&\(536261/280900\)&\(58757/12375\)\\
\bottomrule
\end{longtable}
\endgroup
All \(L\)'s and \(P\)'s are positive.  Each row therefore proves both
the localization of its conditional branch and its Weyl payment.

\section{The connected cap--74 path}

The seventeenth localization row is tied to the exceptional
\(k_9,H=6,h=4\) cap as follows.  Angular propagation from the strict
\((3,2)\) Bessel wall gives
\begin{equation}
 j_{9/2,2}^2>\left(\frac{103}{10}\right)^2+20-12
 =\frac{11409}{100}=108+\frac{609}{100}.
 \label{eq:cap74-ball}
\end{equation}
Accordingly the strict proxy uses
\(\beta_{4,2}^2=11409/100\); the strict inequality belongs to the actual
Bessel zero \(j_{9/2,2}\), not to the proxy constant.
At \(k_9\), the product wall for \((4,2)\) is
\[
 3x<90-20=70,
 \qquad\text{i.e.}\qquad x<\frac{70}{3}.
\]
The ball wall \eqref{eq:cap74-ball}, on the other hand, forces
\[
 x>\frac{11409}{100}-90=\frac{2409}{100},
 \qquad
 \frac{2409}{100}-\frac{70}{3}=\frac{227}{300}>0.
\]
Thus the fully sharpened cap--74 cell is empty.

For completeness, even the deliberately conservative cell is paid.  At
\(\rho=7/20\),
\begin{equation}
 x_{7/20}-\frac{70}{3}
 >\frac{400}{169}\left(\frac{333}{106}\right)^2-\frac{70}{3}
 =\frac{36230}{1424163}>0,
 \label{eq:cap74-rho}
\end{equation}
so a counted \((4,2)\) mode forces \(\rho<7/20\).  Also
\begin{equation}
 108-\left(\frac{207}{20}\right)^2=\frac{351}{400}>0,
 \label{eq:cap74-K}
\end{equation}
so its ball wall forces \(K>207/20\).  Therefore
\begin{align}
 \W(\rho,K)
 &>\frac7{99}\left(1-\left(\frac7{20}\right)^3\right)
       \left(\frac{207}{20}\right)^3\notag\\
 &=\frac{52823261673}{704000000}
 =74+\frac{727261673}{704000000}>74.
 \label{eq:cap74-payment}
\end{align}
The cap itself is exactly
\((6+1)^2+(4+1)^2=49+25=74\).  Equations
\eqref{eq:cap74-ball}--\eqref{eq:cap74-payment} cover both the actual
incompatibility and every row of the conservative cap--74 fallback.

\section{Order-free exhaustion and payment of every high event}
\label{sec:closed-post-event}

This section supplies the finite-combinatorial implication which is not
contained in a list of checkpoint payments alone.  It is deliberately
independent of the order in which two modal walls cross.

For \(7\leq m\leq11\), set
\begin{equation}
 \mathfrak B_m^+
 :=\{(\rho,K):\rho_c\leq\rho\leq39/50,\ 
          k_{m-1}(\rho)\leq K<k_m(\rho)\}.
 \label{eq:post-bands}
\end{equation}
The final spectral band is closed on the right:
\[
 \mathfrak B_{11}^{\rm sp}
 :=\{(\rho,K):\rho_c\leq\rho\leq39/50,\ 
          k_{10}(\rho)\leq K\leq k_{11}(\rho)\}.
\]
At \(K=k_{11}\) we use the original strict proxy.  On the half-open
event bands we use
\[
 b_{\ell,n}(z)^2
 =\max\{n^2z^2+\ell(\ell+1),\,\beta_{\ell,n}^2\},
 \qquad
 C_z(K)=\sum_{\ell,n}(2\ell+1)
             \mathbf1_{\{b_{\ell,n}(z)<K\}},
\]
and define the closed post-event count
\begin{equation}
 \widehat C_z(K)
 :=\sum_{(\ell,n)}(2\ell+1)
       \mathbf 1_{\{b_{\ell,n}(z)\leq K\}}.
 \label{eq:closed-post-count}
\end{equation}
The sum is over the checkpoint inventory appropriate to the upper endpoint
of the band.  Thus a simultaneous event is charged with all of its entering
multiplicities before it is paid.

For a post-event state let \(H\) be the largest first-channel angular index,
let \(h=-1\) if no second channel is present and otherwise let \(h\) be its
largest angular index, and put \(\tau=1\) if a third channel is present and
\(\tau=0\) otherwise.  Missing angular indices only lower the count, since
\begin{equation}
 \widehat C_z(K)
 \leq (H+1)^2+\mathbf1_{\{h\geq0\}}(h+1)^2+4\tau.
 \label{eq:block-post-cap}
\end{equation}

The exhaustive checkpoint inventories are
\begin{equation}
\begin{array}{c|c|c|c}
m&H&h&\text{third channels}\\ \hline
7&H\leq6&h\leq1&\text{none}\\
8&H\leq7&h\leq3&\text{none}\\
9&H\leq8&h\leq4&\text{none}\\
10&H\leq9&h\leq5&\ell=0,1\\
11&H\leq10&h\leq4&\ell=0,1.
\end{array}
\label{eq:event-inventories}
\end{equation}
First channels with \(\ell\geq m\) are excluded at \(k_m\) by their product
wall.  For the first omitted second channel, the two necessary strict bounds
are
\[
\begin{array}{c|ccccc}
m&7&8&9&10&11\\ \hline
\text{first omitted }\ell&2&4&5&6&5\\
\text{product wall}&x<50/3&x<52/3&x<20&x<68/3&x<34\\
\text{ball wall}&x>25&x>4209/100&x>7641/100&
x>6841/100&x>3441/100.
\end{array}
\]
The \(m=10\) ball wall follows from angular propagation of the
\(\ell=5,n=2\) registry wall.  In every column the bounds are incompatible;
for larger \(\ell\), the product upper bound decreases and the propagated
ball lower bound increases.

For \(n=3\), the product wall gives
\(x\leq45/4<(193/53)^2<x\) through \(m=9\), and at \(m=10\) every
\(\ell\geq2\) gives \(x\leq13<(193/53)^2\).  At \(m=11\), the first
omitted third channel \(\ell=2\) would require
\[
 x<63/4,\qquad x>(61/5)^2-132=421/25.
\]
Angular propagation removes the rest of that tail.  For \(n\geq4\), the
least restrictive omitted case is \(m=11,\ell=0,n=4\), which would require
\(x<44/5<(193/53)^2\).  This proves
\eqref{eq:event-inventories} without an unprinted search.

The strict simultaneous restrictions used below are
\begin{equation}
\begin{array}{c|l}
m&\text{restriction}\\ \hline
8&H=6\Longrightarrow h\leq1,\qquad H=7\text{ unrestricted by this row},\\
9&H=8\Longrightarrow h=-1,\qquad H=7\Longrightarrow h\leq2,\\
10&H=9\Longrightarrow h=-1,\tau=0,\quad
   H=8\Longrightarrow h\leq3,\tau=0,\quad H=7\Longrightarrow\tau=0,\\
11&H=10\Longrightarrow h=-1,\tau=0,\qquad H=9\Longrightarrow\tau=0.
\end{array}
\label{eq:event-simultaneous-restrictions}
\end{equation}
For \(m=8,9\) these are exactly the incompatible cells in
\eqref{eq:impossible-cells}.  For \(m=10\), the three first-wall lower
bounds are \(1346/25,1081/36,89/4\), which respectively exceed the
applicable second or third thresholds \(110/3,30,55/4\).  For \(m=11\),
\eqref{eq:H10}--\eqref{eq:H9} give the last two restrictions.  All gaps
are strict, so the same restrictions hold for the closed post-event count
on the half-open bands.

\subsection{The fourteen localization witnesses}

If a closed post-event second mode \(h\) is present below \(k_m\), then
\[
 4x+h(h+1)\leq K^2<x+m(m+1),
 \qquad
 3x<m(m+1)-h(h+1).
\]
At the strict spectral upper face the first inequality is strict and the
same conclusion follows.  Likewise, a third mode forces
\(8x<m(m+1)\).  The following table contains every localization used in
the state table.  Its last column is
\((333/(106(1-r)))^2-A>0\), where \(x<A\) is the forced bound.

\begingroup
\small
\renewcommand{\arraystretch}{1.14}
\begin{longtable}{@{}c L{5.0cm} c c L{3.5cm}@{}}
\toprule
band&state implication&\(A\)&\(r\)&exact positive witness\\
\midrule
\endfirsthead
\toprule
band&state implication&\(A\)&\(r\)&exact positive witness\\
\midrule
\endhead
\(\mathfrak B_7^+\)&some second channel&\(56/3\)&\(3/10\)&\(608779/412923\)\\
\(\mathfrak B_8^+\)&some second channel&\(24\)&\(3/8\)&\(88824/70225\)\\
\(\mathfrak B_8^+\)&\(h\geq2\)&\(22\)&\(1/3\)&\(9233/44944\)\\
\(\mathfrak B_8^+\)&\(h=3\)&\(20\)&\(3/10\)&\(19405/137641\)\\
\(\mathfrak B_9^+\)&some second channel&\(30\)&\(3/7\)&\(40281/179776\)\\
\(\mathfrak B_9^+\)&\(h\geq2\)&\(28\)&\(5/12\)&\(138056/137641\)\\
\(\mathfrak B_9^+\)&\(h\geq3\)&\(26\)&\(2/5\)&\(15889/11236\)\\
\(\mathfrak B_9^+\)&\(h=4\)&\(70/3\)&\(7/20\)&\(36230/1424163\)\\
\(\mathfrak B_{10}^+\)&some second channel&\(110/3\)&\(1/2\)&\(23677/8427\)\\
\(\mathfrak B_{10}^+\)&\(h\geq4\)&\(30\)&\(3/7\)&\(40281/179776\)\\
\(\mathfrak B_{10}^+\)&\(h=5\)&\(80/3\)&\(2/5\)&\(25195/33708\)\\
\(\mathfrak B_{10}^+\)&\(\tau=1\)&\(55/4\)&\(4/25\)&\(65185/275282\)\\
\(\mathfrak B_{11}^+\)&some second channel&\(44\)&\(8/15\)&\(725209/550564\)\\
\(\mathfrak B_{11}^+\)&\(\tau=1\)&\(33/2\)&\(1/4\)&\(5871/5618\)\\
\bottomrule
\end{longtable}
\endgroup
Every defining ratio face is therefore on the excluded side of its state;
the directions in this table are interval implications, not sampled
decisions.

\subsection{Three payment rules}

Put \(F_j(\rho):=\W(\rho,k_j(\rho))\), and define
\begin{equation}
 P(C;c,r):=\frac7{99}(1-r^3)c^3-C,
 \qquad
 Q_j(c,r):=
 \left(\frac{333}{106(1-r)}\right)^2+j(j+1)-c^2.
 \label{eq:PQ-def}
\end{equation}
The derivative calculation in \eqref{eq:Wm-derivative} gives, uniformly
for \(6\leq j\leq10\),
\[
 F_j'(\rho)>0,
 \qquad
 9-\frac{110}{16}=\frac{17}{8}>0,
\]
and \(\W(\rho,c)\) decreases with \(\rho\) at fixed \(c>0\).  Together
with \(2/(9\pi)>7/99\), this gives three rules.

\begin{itemize}
\item[\(L\).] If a state forces \(K\geq c\) and \(\rho<r\), then
      \(P(C;c,r)>0\) implies \(\W(\rho,K)>C\).
\item[\(B\).] If a state forces \(K\geq c\) in \(\mathfrak B_m^+\),
      and \(P(C;c,r)>0\), \(Q_{m-1}(c,r)>0\), then it is paid on both
      sides of \(r\).  For \(\rho\leq r\), use the fixed wall \(c\).
      For \(\rho\geq r\), use \(K\geq k_{m-1}(\rho)\), the increase of
      \(F_{m-1}\), and \(k_{m-1}(r)>c\).  In particular,
      \(F_{m-1}(r)=\W(r,k_{m-1}(r))>\W(r,c)>C\).
      Both arguments are strict at
      \(\rho=r\).
\item[\(G\).] If the state cap is at most a displayed global cap \(C\),
      then \(K\geq k_{m-1}(\rho)\) and the strict global lower-checkpoint
      payment pays it, including the lower checkpoint itself.
\end{itemize}

There is one initial checkpoint anchor of type \(B_0\).  Since
\(\pi<7/2\), one has \(\rho_c>1/8\).  The increase of \(F_6\), together
with the \(B_{7,0}\) row below, gives
\[
 F_6(\rho)>F_6(1/8)>\W(1/8,73/10)>26.
\]

\subsection{The disjoint thirty-eight-state map}

The following rows are literally disjoint.  In a \(G\)-row the displayed
state is the exact complement of the exceptional rows in its band, not a
second overlapping estimate.  The cap column is allowed to dominate more
than one exact block sum, but every exact post-event cap is obtained from
\eqref{eq:block-post-cap} and is one of the block sums represented here.

\begingroup
\scriptsize
\renewcommand{\arraystretch}{1.12}
\begin{longtable}{@{}c c L{6.0cm} c c c@{}}
\toprule
row&band&disjoint post-event state&cap \(C\)&\(c\)&\(r\)\\
\midrule
\endfirsthead
\toprule
row&band&disjoint post-event state&cap \(C\)&\(c\)&\(r\)\\
\midrule
\endhead
\(B_{7,0}\)&7&\(H\leq4,\ -1\leq h\leq0\)&26&\(73/10\)&\(1/8\)\\
\(L_{7,1}\)&7&\(H\leq4,\ h=1\)&29&\(77/10\)&\(3/10\)\\
\(B_{7,1}\)&7&\(H=5,\ h=-1\)&36&\(87/10\)&\(1/2\)\\
\(L_{7,2}\)&7&\(H=5,\ 0\leq h\leq1\)&40&\(87/10\)&\(3/10\)\\
\(B_{7,2}\)&7&\(H=6,\ h=-1\)&49&\(10\)&\(3/5\)\\
\(L_{7,3}\)&7&\(H=6,\ 0\leq h\leq1\)&53&\(10\)&\(3/10\)\\ \midrule
\(G_8\)&8&\(H\leq4,\ h\leq2\), or \(H=5,\ h\leq1\)&40&--&--\\
\(L_{8,1}\)&8&\(H=5,\ h=2\)&45&\(9\)&\(1/3\)\\
\(L_{8,2}\)&8&\(H\leq5,\ h=3\)&52&\(103/10\)&\(3/10\)\\
\(B_{8,1}\)&8&\(H=6,\ h=-1\)&49&\(10\)&\(3/5\)\\
\(L_{8,3}\)&8&\(H=6,\ 0\leq h\leq1\)&53&\(10\)&\(3/8\)\\
\(B_{8,2}\)&8&\(H=7,\ h=-1\)&64&\(23/2\)&\(2/3\)\\
\(L_{8,4}\)&8&\(H=7,\ 0\leq h\leq3\)&80&\(23/2\)&\(3/8\)\\ \midrule
\(G_9\)&9&\(H\leq4\); or \(H=5,h\leq3\); or \(H=6,h\leq1\)&53&--&--\\
\(L_{9,1}\)&9&\(H=6,\ h=2\)&58&\(10\)&\(5/12\)\\
\(L_{9,2}\)&9&\(H=6,\ h=3\)&65&\(103/10\)&\(2/5\)\\
\(L_{9,3}\)&9&\(5\leq H\leq6,\ h=4\)&74&\(207/20\)&\(7/20\)\\
\(B_{9,1}\)&9&\(H=7,\ h=-1\)&64&\(23/2\)&\(2/3\)\\
\(L_{9,4}\)&9&\(H=7,\ 0\leq h\leq2\)&73&\(23/2\)&\(3/7\)\\
\(B_{9,2}\)&9&\(H=8,\ h=-1\)&81&\(71/6\)&\(2/3\)\\ \midrule
\(G_{10}\)&10&\(H\leq4\); or \(H=5\) except \((h,\tau)=(5,1)\); or
\(H=6,h\leq3\); or \(H=7,h\leq2\)&73&--&--\\
\(L_{10,1}\)&10&\(H=6,\ h=4,\ \tau=0\)&74&\(53/5\)&\(3/7\)\\
\(L_{10,2}\)&10&\(H=6,\ h=4,\ \tau=1\)&78&\(53/5\)&\(4/25\)\\
\(L_{10,3}\)&10&\(H=6,\ h=5,\ \tau=0\)&85&\(11\)&\(2/5\)\\
\(L_{10,4}\)&10&\(5\leq H\leq6,\ h=5,\ \tau=1\)&89&\(11\)&\(4/25\)\\
\(L_{10,5}\)&10&\(H=7,\ h=3\)&80&\(23/2\)&\(1/2\)\\
\(L_{10,6}\)&10&\(H=7,\ h=4\)&89&\(23/2\)&\(3/7\)\\
\(L_{10,7}\)&10&\(H=7,\ h=5\)&100&\(23/2\)&\(2/5\)\\
\(B_{10,1}\)&10&\(H=8,\ h=-1\)&81&\(71/6\)&\(3/5\)\\
\(L_{10,8}\)&10&\(H=8,\ 0\leq h\leq3\)&97&\(71/6\)&\(1/2\)\\
\(B_{10,2}\)&10&\(H=9,\ h=-1\)&100&\(64/5\)&\(16/25\)\\ \midrule
\(G_{11}\)&11&\(H\leq6\); or \(H=7\) except \((h,\tau)=(4,1)\); or
\(H=8,\ h\leq1\)&89&--&--\\
\(L_{11,1}\)&11&\(H=7,\ h=4,\ \tau=1\)&93&\(23/2\)&\(1/4\)\\
\(B_{11,1}\)&11&\(H=8,\ 2\leq h\leq4,\ \tau=0\)&106&\(71/6\)&\(3/7\)\\
\(L_{11,2}\)&11&\(H=8,\ 2\leq h\leq4,\ \tau=1\)&110&\(71/6\)&\(1/4\)\\
\(B_{11,2}\)&11&\(H=9,\ h=-1\)&100&\(64/5\)&\(3/5\)\\
\(L_{11,3}\)&11&\(H=9,\ 0\leq h\leq4\)&125&\(64/5\)&\(8/15\)\\
\(B_{11,3}\)&11&\(H=10,\ h=-1\)&121&\(69/5\)&\(13/20\)\\
\bottomrule
\end{longtable}
\endgroup

Apart from the checkpoint anchor \(B_{7,0}\), the wall \(c\) in each
exceptional row is forced by the corresponding fixed part of
\(b_{\ell,n}\).  Besides the displayed first and second walls, the only
weakenings are
\[
 \beta_{4,2}^2-(207/20)^2=2787/400>0,
 \quad
 \beta_{4,2}^2-(53/5)^2=173/100>0,
 \quad
 \beta_{5,2}-11=19/10>0.
\]

\subsection{The twenty-two localized payments}

For each \(L\)-row, the appropriate entry of the fourteen-row table gives
\(\rho<r\), and the fixed wall gives \(K\geq c\).  Direct substitution in
\eqref{eq:PQ-def} gives the following complete list.

\begin{center}
\small
\renewcommand{\arraystretch}{1.12}
\begin{tabular}{@{}c c@{\qquad}c c@{}}
\toprule
row&\(P(C;c,r)\)&row&\(P(C;c,r)\)\\
\midrule
\(L_{7,1}\)&\(21676933/9000000\)&\(L_{7,2}\)&\(58340437/11000000\)\\
\(L_{7,3}\)&\(1564/99\)&\(L_{8,1}\)&\(51/11\)\\
\(L_{8,2}\)&\(2294563597/99000000\)&\(L_{8,3}\)&\(88567/6336\)\\
\(L_{8,4}\)&\(8866645/405504\)&\(L_{9,1}\)&\(162353/21384\)\\
\(L_{9,2}\)&\(10063157/1375000\)&\(L_{9,3}\)&\(727261673/704000000\)\\
\(L_{9,4}\)&\(252947/9702\)&\(L_{10,1}\)&\(2173382/606375\)\\
\(L_{10,2}\)&\(126077081/21484375\)&\(L_{10,3}\)&\(386/125\)\\
\(L_{10,4}\)&\(73838/15625\)&\(L_{10,5}\)&\(89303/6336\)\\
\(L_{10,6}\)&\(97715/9702\)&\(L_{10,7}\)&\(7197/11000\)\\
\(L_{10,8}\)&\(943655/171072\)&\(L_{11,1}\)&\(72407/5632\)\\
\(L_{11,2}\)&\(810599/152064\)&\(L_{11,3}\)&\(32924779/41765625\)\\
\bottomrule
\end{tabular}
\end{center}
All twenty-two fractions are positive, so rule \(L\) pays the dominant
post-event cap in every localized row.

\subsection{The twelve bridge pairs and four global rows}

The eleven ordinary \(B\)-rows and the initial anchor \(B_{7,0}\) have
the following exact \(P,Q\) pairs.  The subscript on \(Q\) is the lower
checkpoint index.

\begingroup
\small
\renewcommand{\arraystretch}{1.12}
\begin{longtable}{@{}c c c@{}}
\toprule
row&\(P(C;c,r)\)&\(Q_j(c,r)\)\\
\midrule
\endfirsthead
\toprule
row&\(P(C;c,r)\)&\(Q_j(c,r)\)\\
\midrule
\endhead
\(B_{7,0}\)&\(73625809/50688000\)&\(Q_6=22025711/13764100\)\\
\(B_{7,1}\)&\(417183/88000\)&\(Q_6=1625379/280900\)\\
\(B_{7,2}\)&\(637/99\)&\(Q_6=165473/44944\)\\
\(B_{8,1}\)&\(637/99\)&\(Q_7=794689/44944\)\\
\(B_{8,2}\)&\(249635/21384\)&\(Q_7=35314/2809\)\\
\(B_{9,1}\)&\(249635/21384\)&\(Q_8=80258/2809\)\\
\(B_{9,2}\)&\(835355/577368\)&\(Q_8=525692/25281\)\\
\(B_{10,1}\)&\(14506973/1336500\)&\(Q_9=4713989/404496\)\\
\(B_{10,2}\)&\(202207748/21484375\)&\(Q_9=648969/280900\)\\
\(B_{11,1}\)&\(507845/261954\)&\(Q_{10}=317585/1617984\)\\
\(B_{11,2}\)&\(25143284/1546875\)&\(Q_{10}=8811001/1123600\)\\
\(B_{11,3}\)&\(151707121/11000000\)&\(Q_{10}=426449/3441025\)\\
\bottomrule
\end{longtable}
\endgroup
Every entry is positive.  Rule \(B\), and the preceding \(B_0\) argument
for the first row, therefore pay all bridge states at the event frequency,
not at a later upper checkpoint.

The four \(G\)-rows use the already displayed global calculation
\eqref{eq:global-coarse}.  For clarity, here are the exact references
needed by the state map; \(S=(193/53)^2+j(j+1)-t^2\) and
\(R=(38/539)t^3-C\).
\begin{equation}
\begin{array}{c|c|c|c|c}
\text{row}&j&C&t&(S,R)\\ \hline
G_8&7&40&208/25&(67049/1755625,\ 5083656/8421875)\\
G_9&8&53&923/100&(1901439/28090000,\ 656778873/269500000)\\
G_{10}&9&73&254/25&(61431/1755625,\ 7911557/8421875)\\
G_{11}&10&89&111/10&(14211/280900,\ 1999489/269500).
\end{array}
\label{eq:event-global-four}
\end{equation}
Both components of every pair are positive.

\begin{proposition}[Closed post-event domination]
\label{prop:closed-post-event-domination}
For every \(7\leq m\leq11\) and every
\((\rho,K)\in\mathfrak B_m^+\),
\[
 \widehat C_z(K)<\W(\rho,K).
\]
At \(K=k_{11}(\rho)\), the original strict proxy also satisfies
\(C_z(K)<\W(\rho,K)\).  Consequently every fixed, moving, or coincident
proxy event in \(k_6\leq K\leq k_{11}\) is paid at its own wall.
\end{proposition}

\begin{proof}
Fix a band.  The upper checkpoint inventory
\eqref{eq:event-inventories} contains every possible event in that band;
the algebraic splits in \eqref{eq:algebraic-splits} cover the two higher
radial tails at their equality faces.  Choose \(H,h,\tau\) only after all
modes at the event have been inserted into \eqref{eq:closed-post-count}.
The simultaneous restrictions
\eqref{eq:event-simultaneous-restrictions} and the mutually exclusive
conditions in the thirty-eight-row table put the resulting state in
exactly one row.  Formula \eqref{eq:block-post-cap} proves its displayed
cap, even if the modal set has gaps.

An \(L\)-row is paid by its localization witness and its positive \(P\).
A \(B\)-row is paid on both closed sides of its split by its positive
\(P,Q\).  A \(G\)-row is paid by \eqref{eq:event-global-four}; the initial
row is paid by \(B_{7,0}\).  These alternatives exhaust the table, proving
the closed post-event inequality.

No ordering premise has entered the argument.  Changing the order of two
fixed or moving walls merely changes which disjoint state is selected.
At a branch equality in the maximum defining \(b_{\ell,n}\), both lower
walls remain valid.  At a ratio face the defining state is impossible by
the strict positive localization witness.  At a bridge face both payment
arguments hold with positive reserve.

At an interior modal wall the original strict proxy assigns equality to
the pre-event cap, whereas the stronger closed post-event cap has just
been paid.  At a lower checkpoint \(K=k_{m-1}\), the closed post-event
version is assigned to the following half-open band and is covered even
when \(K=k_{m-1}\); the original strict proxy omits every channel on its
defining equality and remains owned by the preceding closed spectral
band.  At intermediate upper checkpoints the same convention passes the
post-event state to the following band.  Finally, at \(K=k_{11}\) the
strict inequalities \(b_{\ell,n}<K\) give the same state restrictions and
the \(m=11\) payment rows directly, while the equality channel
\((11,1)\) is absent.  Coincident walls are harmless because
\eqref{eq:closed-post-count} inserts their combined multiplicity before
the unique state row is chosen.
\end{proof}

The proposition above is the final mathematical exhaustion theorem.  It
is not a row-count inference from the historical \(611\)-predicate
executable census.  The map below is retained only as archival provenance
for the older arithmetic implementation and is not a premise of
Proposition~\ref{prop:closed-post-event-domination}.

\section{Archival exact row map and scope audit}

The following map uses the ordering of the two functions in the original
exact registry.  To keep the table readable, \(L\) denotes
\texttt{verify\_localizations}, and \(C\) denotes
\texttt{verify\_cap\_tables\_and\_payments}.  The map lets a referee match
every replaced predicate to a displayed hand calculation.

\begin{center}
\small
\renewcommand{\arraystretch}{1.16}
\begin{tabular}{@{}c c L{12.1cm}@{}}
\toprule
registry family&row numbers&analytic replacement here\\
\midrule
\(L\)&1--2&distinctness and connectedness:
\eqref{eq:ratio-order}--\eqref{eq:ratio-gaps} and the seventeen-row table\\
\(L\)&3--19&all ratios in the high strip, by
\(\rho_c<1/7<4/25\) and \(2/3<7/8\)\\
\(L\)&20--36&the seventeen exact sign witnesses in
the long table\\
\(L\)&37--42&split location
\eqref{eq:split-strip} and side/equality arithmetic
\eqref{eq:algebraic-splits}\\
\(L\)&43--47&the five \(k_{11}\) implications
\eqref{eq:k11-second}--\eqref{eq:H9}\\
\(C\)&44--55&the twelve non-bookkeeping cap--74
predicates, \eqref{eq:cap74-ball}--\eqref{eq:cap74-payment};
row \(C56\), the block identity \(49+25=74\), is allocated to Packet~A\\
\(C\)&57--59&the three impossible
cells, \eqref{eq:impossible-cells}\\
\(C\)&78--79&the general checkpoint
monotonicity proof \eqref{eq:Wm-derivative}\\
\(C\)&80--84&the five global-payment
rows following \eqref{eq:global-coarse}\\
\(C\)&85--98&the seven localization
and seven conditional-payment rows in the final long table\\
\bottomrule
\end{tabular}
\end{center}

The count reconciliation is
\[
\begin{array}{c|c|c}
\text{subfamily}&\text{connected mathematical objects}
&\text{registry predicates}\\ \hline
\text{registry metadata}&1&2\\
\text{high-strip membership}&17&17\\
\text{localization signs}&17&17\\
\text{algebraic splits}&3&6\\
\text{special }k_{11}\text{ implications}&5&5\\
\text{cap--74 path}&1&12\\
\text{impossible cells}&3&3\\
\text{checkpoint monotonicity}&1&2\\
\text{global payments}&5&5\\
\text{conditional payments}&7&14\\ \hline
\text{total}&60&83.
\end{array}
\]
In particular, the five headline Packet--C subfamilies contain exactly
\[
 17+3+3+5+7=35
\]
connected objects (localization signs, algebraic splits, impossible cells,
global payments, and conditional payments).  Because each algebraic split
has a location and a side predicate, and each conditional payment has a
localization and a reserve predicate, these headline objects replace
\(17+6+3+5+14=45\) registry predicates.  The remaining 38 predicates in
the total are the explicitly listed auxiliary closure rows.

Thus this packet replaces all 47 predicates of \(L\) and 36 predicates of
\(C\).  The elementary cap addition rows \(1\)--\(43\), row \(56\), and
rows \(60\)--\(77\) of the latter family are
outside the declared Packet~C scope; they are instances of the general
odd-sum/block-cap identity and are not claimed here.  No row in the
declared scope is unreconstructed.

\end{document}
